% Example APA 7 Manuscript
% Header note: This file is known to compile safely.
\documentclass[12pt,letterpaper]{article}
% APA 7 student manuscript example. Edit the metadata and bracketed content.
% Ready to build: Example_References.bib is connected below.
% Run ./compile_manuscript.sh or PDFLaTeX/Biber/PDFLaTeX/PDFLaTeX.
% This file configures a standard article class explicitly for student papers.
\usepackage[T1]{fontenc}
\usepackage[utf8]{inputenc}
\usepackage{mathptmx} % 12-point Times-family text; PDFLaTeX-compatible.
\usepackage[margin=1in,headheight=15pt,headsep=0.25in]{geometry}
\usepackage{setspace}
\usepackage{indentfirst}
\usepackage{fancyhdr}
\usepackage{titlesec}
\usepackage{booktabs,array}
\usepackage{graphicx,float}
\usepackage{listings}
% APA 7 bibliography and example reference file.
\usepackage[american]{babel}
\usepackage{csquotes}
\usepackage[backend=biber,style=apa]{biblatex}
\DeclareLanguageMapping{american}{american-apa}
\addbibresource{Example_References.bib}
\usepackage[hidelinks]{hyperref}
\urlstyle{same}

% Student-paper header: page number at upper right.
\pagestyle{fancy}
\fancyhf{}
\fancyhead[R]{\thepage}
\renewcommand{\headrulewidth}{0pt}
\renewcommand{\footrulewidth}{0pt}
\doublespacing
\raggedright
\setlength{\parindent}{0.5in}
\setlength{\parskip}{0pt}
\setcounter{secnumdepth}{0}
% APA heading levels 1--3.
\titleformat{\section}{\normalfont\normalsize\bfseries\filcenter}{}{0pt}{}
\titleformat{\subsection}{\normalfont\normalsize\bfseries}{}{0pt}{}
\titleformat{\subsubsection}{\normalfont\normalsize\bfseries\itshape}{}{0pt}{}
\titlespacing{\section}{0pt}{0pt}{0pt}
\titlespacing{\subsection}{0pt}{0pt}{0pt}
\titlespacing{\subsubsection}{0pt}{0pt}{0pt}
\setlength{\bibhang}{0.5in}
\setlength{\bibitemsep}{0pt}
\AtBeginBibliography{\doublespacing}
\lstset{language=bash,basicstyle=\ttfamily\footnotesize,
  breaklines=true,columns=fullflexible,keepspaces=true,
  showstringspaces=false,frame=none,aboveskip=0.5\baselineskip,
  belowskip=0.5\baselineskip}

% Student and course details.
\newcommand{\papertitle}{Example APA 7 Manuscript}
\newcommand{\studentname}{[Student Name]}
\newcommand{\affiliation}{[Department or Program], [Institution]}
\newcommand{\coursename}{[Course Name]}
\newcommand{\instructorname}{[Instructor Name]}
\newcommand{\duedate}{[Month Day, Year]}
% Abstracts are not normally required for student papers unless assigned.
\newif\ifincludeabstract
\includeabstractfalse % Change to \includeabstracttrue if required.

\begin{document}

% Title page: bold centered title, one blank double-spaced line before author.
\begin{center}
\vspace*{2\baselineskip}
\textbf{\papertitle}\par
\vspace{\baselineskip}
\studentname\par
\affiliation\par
\coursename\par
\instructorname\par
\duedate\par
\end{center}
\newpage

\ifincludeabstract
\section*{Abstract}
\noindent [Write one unindented paragraph summarizing the purpose, design,
data, methods, principal results, and conclusion. Use the word limit assigned
by the instructor. Replace this instruction only after verifying the results.]
\par\noindent\textit{Keywords:} [term one], [term two], [term three]
\newpage
\fi

% Begin the text under the repeated paper title.
\section*{\papertitle}
This template organizes a student manuscript around an independently chosen
research question and permitted data. Replace the bracketed guidance
with your question, methods, verified results, and interpretation. The
Dr. Seuss references illustrate narrative citation, \textcite{seuss1957cat},
and parenthetical citation \parencite{seuss1971lorax}. Use sources relevant
to your study when preparing your manuscript.

[Introduce the problem, explain why it matters, and summarize relevant prior
work with citations. End this opening section with the research question or
objective. For this independent example, the objective is to inspect a permitted sample table,
select records using a documented rule, and verify the resulting files.]

\section{Method}
\subsection{Design and Data Source}
[Describe an independently selected study or demonstration dataset. Identify
its public source or your right to use it, its unit of observation, field
schema, collection process, and limitations. Do not insert restricted course
datasets or original assignment requirements into a public template.]

[Report acquisition date, inclusion criteria, sample size, and relevant
permissions. For an empirical study, describe the sampling design and ethics
approval or exemption when applicable.]

\subsection{Measures and Variables}
[Define the variables relevant to the question, their units, category labels,
missing-value conventions, and any derived measures. Distinguish identifiers
from measurements. State the source of definitions rather than inventing units.]

\subsection{Procedure}
[Document the steps used to validate a permitted source file, inspect its
schema, apply justified inclusion rules, and check outputs. Use independently
written examples rather than reproduced classroom commands or solutions.]

[Describe the exact software versions, commands or functions, file locations,
validation checks, and deviations from the planned procedure. Simple AWK and
CUT processing is suitable only for the documented simple CSV layout; a
CSV-aware parser is needed for embedded commas or multiline fields.]

\subsection{Analysis Plan}
[State the descriptive summaries or statistical models used, the treatment of
missing data, assumption checks, significance threshold if applicable, and
multiple-testing adjustment if relevant. Explain the design rather than
listing software packages. A descriptive file-processing example may verify records and files without
requiring a hypothesis test.]

\section{Results}
[Report only output that was actually produced and checked. Begin with the
number of source and eligible records, then present the findings in the order
of the research question. Refer to each table and figure in the text. Do not
repeat every table entry in prose.]

Table~\ref{tab:validation} illustrates a results-table layout. Its cells are
placeholders, not measured values. Figure~\ref{fig:example} illustrates the
position of a figure number, title, image, and note. Replace the image area
with a verified plot before submission, or remove the figure if unnecessary.

% Table/figure bodies may be single-spaced for readability; main text stays double-spaced.
% APA: bold number, italic title above; explanatory note below. No vertical rules.
\begin{table}[H]
\refstepcounter{table}\label{tab:validation}
\noindent\textbf{Table \thetable}\\
\noindent\textit{Metadata Processing and Validation Summary}\par
\begin{singlespace}
\smallskip
\begin{tabular}{@{}p{0.42\linewidth}p{0.53\linewidth}@{}}
\toprule
Validation item & Observed result \\
\midrule
Source records, excluding header & [Insert verified count] \\
Selected records, excluding header & [Insert verified count] \\
Header retained & [Insert check outcome] \\
Condition and age extracted & [Insert check outcome] \\
Subset files retained after cleanup & [Insert check outcome] \\
\bottomrule
\end{tabular}
\smallskip
\end{singlespace}
\noindent\textit{Note.} Record counts exclude the header; a newline count may
include it. Replace every cell after inspecting the actual outputs. The table
describes processing checks, not population estimates.
\end{table}

\begin{figure}[H]
\refstepcounter{figure}\label{fig:example}
\noindent\textbf{Figure \thefigure}\\
\noindent\textit{Example Figure Area for a Verified Analysis Plot}\par
\begin{singlespace}
\smallskip
\noindent\fbox{\parbox[c][1.15in][c]{0.94\linewidth}{\centering
Replace this area with a verified figure.\\
No data or results are displayed here.}}
% Replace the fbox above with, for example:
% \includegraphics[width=\linewidth]{figures/your_verified_plot.pdf}
\smallskip
\end{singlespace}
\noindent\textit{Note.} Describe the observations, axes, units, groups, scale,
and abbreviations needed to read the figure. State whether the data are
simulated. Cite an adapted or reproduced figure and check its permission terms.
\end{figure}

\section{Discussion}
\subsection{Principal Findings}
[Answer the research question using the verified results. Explain the meaning
of the findings, compare them with prior work using appropriate citations,
and distinguish a descriptive pattern from inferential evidence.]

\subsection{Limitations}
[Explain limitations of the design, input quality, missing data, assumptions,
and generalizability. When using synthetic records, avoid claims about real populations. Simple
delimiter tools do not implement general CSV parsing. Explain how these constraints affect the interpretation.]

\subsection{Implications and Conclusion}
[Explain the practical or methodological implication and the next justified
step. Close with a concise answer to the objective; avoid introducing new
results or making claims beyond the study design.]

% References: new page, bold centered heading, alphabetical entries,
% double spacing and 0.5-inch hanging indents. Cite every retained entry.
\newpage
\section*{References}
\printbibliography[heading=none]

% Appendix: separate page after references. Code is supplementary material.
\newpage
\section*{Appendix}
\section*{Reproducible Command Example}
These generic inspection commands assume a separately supplied, permitted
sample file under the ignored \texttt{data/} directory. They neither execute
nor reproduce coursework.

\begin{singlespace}
\begin{lstlisting}
# Independently prepared example commands (adapt locally).
# The CSV is a user-supplied, permitted input outside this repository.
head -n 5 data/example_table.csv
wc -l data/example_table.csv
# Document any filtering or validation separately in your private workspace.
\end{lstlisting}
\end{singlespace}

[Record the operating system, Bash version, source-file checksum, execution
date, and validation outcomes. Include this appendix when your independently prepared manuscript
requires supplementary commands.]

\end{document}
